\documentclass[10pt]{article}
\usepackage[margin=20mm]{geometry}
\usepackage{amsmath,amssymb,fontspec,unicode-math,xurl,hyperref}
\IfFontExistsTF{DejaVu Serif}{\setmainfont{DejaVu Serif}}{\setmainfont{Latin Modern Roman}}
\setmathfont{latinmodern-math.otf}
\setlength{\parindent}{0pt}\setlength{\parskip}{6pt}\sloppy
\title{Technical Verification of Complex Sketches for Dynamic-Weight Distances}
\author{AI-assisted technical review}
\date{}
\begin{document}
\maketitle
Abstract: The checked fourth-root estimator is unbiased, and its released scalar implementation matches the main formula. A substantive limitation arises in the improved block method: balance of the input points does not imply balance of their differences. A fully admissible pair of balanced points has relative distance error at least one with probability one under dimensions satisfying the advertised improved scaling. This invalidates the unrestricted pairwise extension, assessed Major, while leaving the balanced-single-vector theorem intact. A stronger main-text tail formula is already corrected in the appendix and is a local Minor inconsistency. Released experimental settings also differ from the documented settings, with historical figure provenance unresolved. Adjudication remains pending.\par
\section*{Identity and scope}
Simone Moretti, Paolo Pellizzoni, and Francesco Silvestri, Dimensionality Reduction on Complex Vector Spaces for Euclidean Distance with Dynamic Weights, ICML 2025, PMLR 267 [paper]. Assignment 0215 / I25-43584; OpenReview xKMMGugUgy. The reviewed publisher PDF, acquired 25 October 2026 UTC, has 20 pages, including appendices on pp. 12–20. Page numbers below are one-based PDF pages and equal the printed local numbers.\par
The released repository `aidaLabDEI/complex-dimensionality-reduction` is frozen at commit `70cb2255c08ae29f22c2372225f649c90fc5da04` (28 May 2025). Its complete acquired snapshot contains two files, including the 141-line . The independent check read this source without executing it. Source archives and member/version mappings remain unchanged. The acquisition recorded HTTP 403 for separate OpenReview material; no assertion of its absence is made. Later revisions and historical experiment checkpoints were not retrieved.\par
The review checks unbiasedness, the block estimator, assumptions and application of the near-uniform theorem, selected concentration dependencies, finite-family probability accounting, and experimental source settings. It does not certify every imported sub-Weibull result or every coefficient expansion in the general chaos analysis. No author programs, training runs or dataset downloads were performed.\par
\section*{Estimator and necessary assumptions}
Inputs and weights are real, with nonnegative weights in the intended weighted-distance setting: \(\|x\|_w^2=\sum_jx_j^2w_j^2\). The sketch is complex. For independent uniform fourth roots \(A_{ij}\in\{1,-1,i,-i\}\), \[g(x)=Ax/\sqrt{k},\qquad
 \rho(g(x),w)=\frac1k\sum_{r=1}^k
 \operatorname{Re}\left[
       \left(\sum_jA_{rj}x_j\right)^2
       \left(\sum_j A_{rj}^2w_j^2\right)\right].\] Ordinary squares, not Hermitian squared magnitudes, are intentional. Since \(\mathbb EA=\mathbb EA^2=\mathbb EA^3=0\) and \(\mathbb EA^4=1\), only equal-index terms survive expectation. Hence \(\mathbb E\rho=\sum_jx_j^2w_j^2\). Independent enumeration for \(x=(1,2,-1)\) and \(w=(2,1,3)\) over all \(4^3=64\) root triples returned mean \(17\), equal to the exact target. This is a finite corroboration of the moment proof.\par
The block method partitions coordinates into \(L\) fixed blocks, applies independent \(k\)-row sketches, and sums their norm estimates. Section 4.4 assumes a fixed constant \(\zeta\) with \[\|x^{(\ell)}\|_2^2\le\zeta\|x\|_2^2/L,\qquad
 \|w^{(\ell)}\|_4^4\le\zeta\|w\|_4^4/L
 \quad(\ell=1,\ldots,L).\] Theorem 4.6 uses, with \(h=\log(8/\delta)\), \[k\ge c\max\left\{\frac{\Delta^2h}{L^2\epsilon^2},
                  \frac{\Delta h^2}{L^{3/2}\epsilon}\right\}.\] It yields additive error \(\epsilon\|x\|_2^2\|w\|_4^2/\Delta\), and relative error when the corresponding norm ratio is at most \(\Delta\). Constants may absorb the fixed balance constant. The theorem is read with the immediately preceding balance assumptions; omitting those assumptions from its repeated statement does not license applying it to arbitrary differences.\par
\section*{Substantive findings}
Problem 1 (Balanced points need not have balanced differences).  The pairwise-distance remark immediately after Theorem 4.6, p. 8, extends the improved guarantee using only a common partition and linearity. Those conditions are insufficient.\par
\subsection*{Source claim and fully specified counterexample}
The remark asserts that fixing the partition for all \(x\in X\) preserves linearity and that the dimensionality reduction maintains pairwise distances in \(X\). Unlike the general estimator, the improved theorem has an additional input-balance hypothesis. Linearity gives \(\bar g(x)-\bar g(y)=\bar g(x-y)\), but does not transfer that hypothesis.\par
Let \(L\) be a positive integer, \(d=4L\), and partition coordinates into \(L\) blocks of size four. Set \[k=1,\qquad w=y=\mathbf1_d,\qquad x=\mathbf1_d+e_1,
 \qquad X=\{x,y\}.\] All weights are positive and the partition is common. The \(y\) and \(w\) block bounds hold with \(\zeta=1\). For \(x\), \(\|x\|_2^2=4L+3\); its first block has squared norm seven and the others four. With \(\zeta=2\), \[7\le2(4L+3)/L,\qquad 4\le2(4L+3)/L .\] Thus both original points satisfy (2.1) with a fixed constant independent of \(L\).\par
Their difference is \(u=x-y=e_1\), with \(\|u\|_w^2=1\) and \(\|u\|_2^2\|w\|_4^2=2\sqrt L\). Choose \(\Delta=2\sqrt L\), any fixed \(0<\epsilon<1\), and, for example, \(\delta=1/4\). Then (2.2) at \(k=1\) is satisfied whenever \[L\ge c\max\{4h/\epsilon^2,\;2h^2/\epsilon\}.\] For any fixed admissible universal constant \(c\), such integers \(L\) exist. The output has \(L\) complex coordinates, or \(2L\) real coordinates, fewer than the input’s \(4L\) real coordinates.\par
Only the first block contributes to the distance estimate: \[\begin{aligned}
 \bar\rho(\bar g(x)-\bar g(y),w)
 &=\operatorname{Re}\{A_1^2(A_1^2+A_2^2+A_3^2+A_4^2)\}\\
 &=1+\sum_{j=2}^4 A_1^2A_j^2 .\end{aligned}\] The three products are independent uniform signs (condition on \(A_1^2\)). The estimate therefore takes values \[-2,\ 0,\ 2,\ 4
 \quad\text{with probabilities}\quad
 \tfrac18,\ \tfrac38,\ \tfrac38,\ \tfrac18 .\] Every value is at distance at least one from the true squared distance one. Consequently the claimed relative-error event fails with probability one for \(\epsilon<1\), regardless of how large \(L\) becomes. Parent enumeration of all 256 root quadruples returned counts \(32,96,96,32\) and the same failure probability.\par
\subsection*{Impact, counterevidence, and exact scope repair}
This is a Major limitation of the improved pairwise conclusion, a central advertised use of linear dimension reduction. It is not a counterexample to Theorem 4.6 applied to a genuinely balanced single vector: \(u=e_1\) has all its energy in one block and fails (2.1) for any fixed \(\zeta\) as \(L\) grows. The general single-vector theorem and its general pairwise application are also not refuted.\par
A correct pairwise extension requires (2.1) for every relevant nonzero difference \(x-y\), and for each relevant weight vector, together with the norm-ratio bound for those differences. Apply Theorem 4.6 to each such fixed pair and use a union bound for a finite collection. This verifies the extension under strengthened assumptions. Reading those difference assumptions implicitly into the remark gives a valid interpretation; balance of the points alone does not. The appendix experiments estimate single-vector norms and do not test this missing pairwise condition.\par
Problem 2 (The main \(L^3\) tail display conflicts with the appendix’s correct \(L^2\) version).  Lemma 4.5 has different Gaussian exponents on pp. 8 and 18.\par
Let \(M=\|x\|_2^2\|w\|_4^2\). The p. 8 display has Gaussian term \(8\exp(-kc_1L^3t^2/M^2)\), whereas the p. 18 restatement has \(8\exp(-kc_1L^2t^2/M^2)\). Both retain the heavy-tail term \(8\exp(-\sqrt{k}c_2 L^{3/4}\sqrt t/\sqrt M)\).\par
The stronger main display is not generally valid. Use \(d=2L\), blocks of size two, \(k=1\), and \(x=w=\mathbf1_d\), which satisfy exact block balance. One block estimate is eight with probability \(1/4\) and zero with probability \(3/4\). Its centered value \(Z\) is six or minus two, so \(\mathbb EZ^2=12\) and \(\mathbb EZ^4=336\). For independent block errors \(S=\sum_{\ell=1}^L Z_\ell\), \[\mathbb ES^2=12L,\qquad \mathbb ES^4=432L^2-96L .\] Applying Cauchy–Schwarz to the truncated second moment, equivalently the Paley–Zygmund inequality for \(S^2\), yields \[\Pr(|S|>\sqrt{6L})\ge
 \frac{(1-1/2)^2(12L)^2}{432L^2-96L}>\frac1{12}.\] But \(M=2\sqrt2L^{3/2}\); at \(t=\sqrt{6L}\) both terms in the p. 8 maximum tend to zero for any fixed positive \(c_1,c_2\). Its purported upper bound contradicts this fixed lower bound.\par
The supplied appendix itself resolves the discrepancy. Lemma A.5 gives fluctuation scale \[M\left(\frac{c_1\sqrt v}{\sqrt{k}L}
       +\frac{c_2v^2}{kL^{3/2}}\right);\] Lemma 4.5 on p. 18 uses \(L^2\), and the final proof on p. 19 explicitly uses that \(L^2\) exponent. Theorem 4.6’s stated dimension requirement agrees with the appendix. Thus this is Minor as a local inconsistent display: the checked operative derivation already uses the correct version. It is not treated as a failure of the final balanced-vector rate, nor as an uncredited new repair.\par
Problem 3 (Released experimental settings differ from the stated protocol).  The available source does not reproduce several declared settings without changes.\par
Appendix B, p. 19, and the figure captions specify 1,000 independent sketch trials. At the pinned commit, `complex\_jl.py:46` sets `N = 50`. Figure 5, p. 20, specifies \(m=100\) nonzero coordinates; line 121 sets `m = 200`. These are direct source/protocol discrepancies. The common dimension \(d=20000\) and the estimator/normalization match the paper.\par
There is also an internal weight-range discrepancy: Appendix B prose states \([0.1,1.2]\), whereas the captions and source line 48 use \([1.0,1.2]\). Both cannot be the exact specification of the same draw. The latter range is supported by the released generator, but the historical draw is not identified from that observation.\par
These are Minor reproducibility discrepancies, with no demonstrated failure of the theoretical estimates or core empirical conclusion. The inspected snapshot is complete but contains no plotting script or hash-bound results identifying which settings generated the published figures. It may represent a cheaper test configuration. Matching `N` and `m` would align those two settings, not regenerate the original plots or settle the weight-range provenance. The saved `out[0]` metadata is an unsquared norm, while samples estimate squared norms; without plotting code, a wrong plotted reference line is not established and remains Needs Evidence.\par
\section*{Other checked scope and remaining work}
The main estimator’s real-part quartic expression and the block row averaging agree with the released code. In particular, replacing squares by complex conjugates would alter the moment construction; no such correction is justified. The new concentration results are conditional on fixed inputs and independent sketch randomness. The paper’s finite-weight-set corollary does not assert validity for unlimited adaptively chosen weights; this review makes no such accusation.\par
For precise finite-family probability bookkeeping, let \(N\) bound the number of pair/weight events and let \(\eta\) be the desired total failure probability. Applying the pointwise theorem at \(\delta=\eta/N\) gives success at least \(1-\eta\). Corollary 4.4’s printed choice \(\delta=1/(|W|n^2)\), p. 7, only gives a failure upper bound of one after the stated crude union bound, not a specified high-probability level. The exact choice above restores a meaningful bound. If \(|W|=O(n)\) and \(\eta=n^{-a}\) for fixed \(a>0\), the logarithmic dimension order remains unchanged. The asymptotic corollary is therefore supported by the already available pointwise result with explicit probability allocation; no false corollary is claimed.\par
No independent empirical distributions were generated using author code. The finite enumerations here test algebra and furnish counterexamples, not the paper’s large-dimensional experiments. A broader test of pairwise distances would require a dataset whose differences satisfy the strengthened assumption, or a different analysis for unbalanced differences.\par
Object | Assessment | Repair verification and limits\par
Improved pairwise extension | Major | Valid with balance of all differences; point balance alone insufficient.\par
Main \(L^3\) display | Minor | Correct \(L^2\) version already supplied and used in appendix.\par
Released protocol settings | Minor | Exact differences verified; historical figure settings unresolved.\par
General estimator | Scoped NIF | Moment identity and static implementation agree; no full chaos-proof certification.\par
All adjudications remain pending; human reviews are empty and author adoption has not been requested. No assertion of first discovery or external erratum status is made.\par
\section*{Supporting checks}
The source file has Git blob `98794c35892c09967ba4693f81f84c2e789a6aea`. Actual staged bytes, archive members and member maps were checked; no source mutation occurred.\par
An independent calculation exhaustively enumerates root tuples, records estimator distributions and checks the moment calculations. Source formula images of pp. 8 and 18 were inspected directly.\par
No author code was executed and no additional source was downloaded.\par
Full role and prompt/source bindings are in . This is an initial review, without acceptance or publication.\par
9 S. Moretti, P. Pellizzoni, and F. Silvestri, Dimensionality Reduction on Complex Vector Spaces for Euclidean Distance with Dynamic Weights, PMLR 267 (2025). Official publisher PDF, frozen retrieval 25 October 2026. https://proceedings.mlr.press/v267/moretti25a.html. aidaLabDEI, complex-dimensionality-reduction, commit `70cb2255c08ae29f22c2372225f649c90fc5da04`, 28 May 2025. https://github.com/aidaLabDEI/complex-dimensionality-reduction.\par
\end{document}
